\section{Introduction}
An operating contract is evaluated through the decisions it induces. A financial guarantee can discourage surrender while reducing the agent's private value; a compulsory operating term can replace a guarantee while imposing a real enforcement cost. When consumption, portfolio choice, preference adjustment, and surrender are jointly chosen, a comparison of private values does not establish a comparison of purchased service. Moreover, an accurate approximation to an optimal policy need not approximate the service of every optimal policy. At an indifference point, an arbitrarily small value error can conceal an economically important change in output.

This paper develops a response-set approach to this problem. The object passed from the dynamic program to the purchaser is not a selected policy or a numerical derivative. It is a set of jointly feasible service and surrender moments constrained by certified value queries. A finite family of such queries gives a linear description, in an enlarged space, of every response consistent with the value information. Its support bounds can be checked by rational weak duality. The purchaser then compares the least favorable response to an executable offer with the most favorable response to each competing offer, including nonprocurement. The same construction admits continuous enforcement instruments: interval bounds cover the instruments between evaluated contracts rather than treating the evaluation grid as the feasible set.

There are three parts to the argument. First, Theorem~\ref{thm:oracle} characterizes the sharp response set relative to a finite value oracle for affine reward perturbations. The characterization accommodates ties, multidimensional service, and an explicit tolerance for near-optimal behavior. Its sharpness is an information statement: every point in the set is generated by an affine-response economy consistent with the queries. Additional knowledge of a particular dynamic economy can, and in our application does, tighten this set. Second, Theorems~\ref{thm:procurement} and~\ref{thm:continuous} convert these response restrictions into procurement and continuous-instrument certificates. An enforcement-capacity result separates the charge imposed on surrender from the capacity required to support it. Third, the numerical application evaluates these comparisons in one stochastic settlement economy, with an unchanged collection of rewards, transitions, financial interpolation rules, and proposal actions.

The application makes the distinction between institutional restrictions and economic primitives consequential. On the original six-point enforcement menu, the optimal compulsory terms are three intervals with preference adjustment and four without it. Refining the charge menu to increments of 0.05 changes this comparison to two versus four. Extending the charge to a continuum changes the optimization problem again. We report an executable contract, an upper bound over every charge in the continuum, and a separate upper bound for each compulsory term. Thus the coarse-menu result is retained as a theorem about its stated institution, while the continuous calculation tests, rather than assumes, its robustness. Neither a separate analytical example nor a favorable local perturbation substitutes for that test.

The service measures also enter the result substantively. Discounted operating duration is the output bought in the baseline contract. It is observed as continued operation, not inferred from the agent's reported utility. Its coefficient in the agent's flow payoff supplies a computational instrument for identifying its response. Discounted elective surrender and wealth-qualified duration are additional moments. A purchaser who receives the surrender fee values duration and surrender jointly. Queries in the joint direction yield strictly tighter bounds than coordinate-wise restrictions alone; the settlement experiment demonstrates this reduction within the actual application. Changes in fee incidence, capacity-cost curvature, compulsory-term costs, outside value, and output quality then reveal which parts of the procurement comparison are institutional and which are incentive effects.

Neural Bellman Operators provide the proposal, evaluation, improvement, and certification architecture producing the value information. The neural component and the certificate have different jobs. A proposal can reduce the search for good feasible lower policies. It cannot establish its own approximation error or the infeasibility of a better response. The verification layer therefore remains valid for a poorly trained network and for a nonneural proposal. This separation preserves the operator and approximation results developed in the accompanying compendium. It also permits an informative unfavorable benchmark: on the deposited settlement target, the fitted network is less accurate than the nearest-anchor policy, and exact upper-value calculation dominates their online comparison. We report these costs rather than attributing the certificate's validity to neural prediction accuracy.

The adjustment mechanism is examined separately from the purchaser's optimization. A signed opportunity-exposure identity separates the current adjustment gain, exposure to future adjustment opportunities, and the replacement of the financial action. Removing downward drift does not by itself rank the two initial financial mandates. We provide primitive sufficient conditions for an ordered opportunity process in a specified stopping subclass, and an explicit 54-point map in the settlement economy, including reversals. The map is not promoted into a uniform theorem between its points. Its adverse cases identify the limits of a mechanism that would otherwise be inferred from a single favorable calibration.

Finally, numerical error must be distinguished from model error. We integrate the data-dependent binary64 error derivation and the strictly positive initial-policy construction into the statements used by the paper. We also give a model-to-decision transfer theorem with separate entries for time, state, action, boundary, terminal, and arithmetic errors. A direct moment audit of the stored transition exposes a material projection effect. It prevents the error for evaluating fixed arrays from being mistaken for the error of approximating the antecedent diffusion. The continuous-state theory is retained, but a numerical conclusion about that diffusion requires the corresponding approximation entries, not a relabeling of the finite-state certificate.

\paragraph*{Related literature.}
The affine dependence of rewards on service moments connects the analysis to reward transfer and successor features; see \citet{alegre2022}. The envelope theorem of \citet{milgrom2002} applies to arbitrary choice sets at differentiability points and supplies a broader foundation for comparative statics. Our finite-query result instead asks what can be established without differentiability from finitely many interval observations, and retains the joint response set needed by a purchaser. Regional verification of parametric Markov models is an established subject. Parameter lifting replaces fixed but unknown transition parameters by state-wise nondeterministic choices; \citet{quatmann2016} develop this approach, and \citet{junges2019} study coverage of parameter regions. We compare an implementation of that relaxation with the common-law count and signed-chord bounds; we do not describe our implementation as an execution of their software. The full discussion of neural policy iteration, nonlinear operator approximation, recursive utility, and dynamic equilibrium is retained in the complete historical compendium.

\paragraph*{Organization.}
Section~\ref{sec:model} specifies the institution and the settlement economy. Section~\ref{sec:response} establishes response identification and procurement. Section~\ref{sec:continuous} treats continuous enforcement. Section~\ref{sec:results} presents the economic evidence. Sections~\ref{sec:mechanism} and~\ref{sec:computation} examine the mechanism and computational architecture. Section~\ref{sec:transfer} gives the model-to-decision error account. The current supplement supplies the proofs and verification details; the historical compendium and supplements retain all earlier material without silently changing its target identity.

\section{The Economic Environment}
\label{sec:model}
\subsection{Operating decisions and affine response moments}
Consider a finite-horizon controlled process with live state $x$ at dates $n=0,\ldots,N$, terminal payoff $G$, and feasible history-dependent operating and stopping policies $p\in\mathcal P_j^r$. The public regime $r$ specifies whether deliberate preference adjustment is available. The instrument $j$ includes the initial financial mandate and the first date on which elective surrender is permitted. It does not specify the agent's subsequent consumption, financial position, or stopping rule. Policies may randomize. The finite-state application below has finite later-date action menus and a continuous, piecewise-affine first-date financial operator.

For a fixed feasible policy, let $A(p)$ be expected discounted operating duration, $H(p)$ discounted elective-surrender incidence, and $Q(p)$ discounted duration during which a specified observable quality criterion is met. These quantities satisfy $0\leq A,H\leq1$ and $0\leq Q\leq A$. Forced discharge is not elective surrender and does not enter $H$. At operating-benefit parameter $d$ and surrender charge $F$, private payoff is
\begin{equation}
 J_j^r(p;d,F)=B_j^r(p)+dA_j^r(p)-FH_j^r(p),\qquad
 W_j^r(d,F)=\sup_{p\in\mathcal P_j^r}J_j^r(p;d,F).
 \label{eq:payoff}
\end{equation}
The feasible policies and their laws are unchanged when only $(d,F)$ is perturbed. This qualification is essential: a perturbation that changes feasible actions, wealth dynamics, or the stopping technology does not identify service through the same affine equation.

More generally, write $J(p;t)=\alpha(p)+t\cdot z(p)$, where $t\in\R^k$ is a reward perturbation and $z(p)\in K$, a known compact polytope. The point $t=0$ denotes the contract under consideration. For duration and surrender, $t=(\Delta d,-\Delta F)$ and $z=(A,H)$. A quality premium can be included by appending its reward feature and perturbation. The results below require neither a smooth optimizer nor a differentiable value function.

\subsection{Enforcement, transfers, and participation}
The purchaser chooses a mandate $\sigma$, compulsory term $m$, charge $F$, enforcement capacity $E$, and signing transfer $q\geq0$, or does not procure. The enforcement technology is
\begin{equation}
 0\leq F\leq\zeta E,\qquad F\leq\bar F,\quad E\geq0,
 \label{eq:enforcement}
\end{equation}
where $\zeta>0$ measures effective charge capacity and $\bar F$ is the administrative charge ceiling. The baseline and incidence counterfactuals use $\bar F=1$; capacity itself can exceed one when enforcement is inefficient. Capacity costs $C(E)$, a nondecreasing function; compulsory terms cost $D(m)$. A fraction $\alpha_C\in[0,1]$ of capacity cost is borne by the agent in a separate utility numeraire, and the remainder by the purchaser. The purchaser receives fraction $\rho_F\in[0,1]$ of elective surrender fees. The agent's outside value is $G_0$.

The signing transfer and capacity expenditure do not alter managed wealth or the operating feasible set in the baseline institution. With an attained exact best response, the infimum transfer inducing weak participation is
\begin{equation}
 q_j=[G_0-W_j(d,F)+\alpha_C C(E)]_+.
 \label{eq:participation}
\end{equation}
Weak participation is resolved by acceptance. A strict participation premium must be added when that convention is inappropriate; our surplus margins provide the budget for such a premium. The purchaser's payoff under response $p$ is
\begin{equation}
 \Pi_j(p,q)=bA(p)+\rho_F FH(p)+\nu Q(p)-q-(1-\alpha_C)C(E)-D(m),
 \label{eq:buyer}
\end{equation}
with $b\geq0$ and a stated quality valuation $\nu$. All comparisons retain the positive part in \eqref{eq:participation}. A transfer-timing change is neutral only when the discount factors and numeraire conversion are specified; a payment that enters managed wealth is a different dynamic problem and must be re-solved.

Equation~\eqref{eq:buyer} is not an unrestricted mechanism-design formulation. The purchaser does not elicit a private type, select a contingent reporting mechanism, or contract directly on preference adjustment. It is a complete-information choice of enforceable operating instruments with endogenous response. Its economic content comes from comparing these instruments before selecting an allocation and from respecting their costs and participation effects.

\subsection{The stochastic settlement specification}
The quantitative economy uses $N=8$ intervals of length $1/8$, state $(u,X)$, and operating action $(c,\theta,\pi)$. The preference state $u$ lies on 33 nodes in $[1.2,2.8]$ and wealth on 49 nodes in $[0.5,2]$, giving 1,617 states. The initial state is $(2,1.25)$. The cardinal operating primitive and terminal payoff are
\begin{equation}
 f(c,u,\theta;d)=\frac{c^{1-u}}{1-u}-\frac{k}{2}\theta^2+d,
 \qquad G(u,X)=-0.02(u-2)^2+0.1\log X.
 \label{eq:primitives}
\end{equation}
The adjustment-cost coefficient is $k=2$. The canonical constructor fixes discounting, drift, volatility, and discharge rewards. The rate of time preference is $0.04$; the financial drift parameters are $0.02$ and $0.06$, and volatility scales are $0.05$ and $0.2$ for preferences and wealth. The two endpoint correlation laws are $-0.25$ and $0.25$.

For endpoint $i\in\{0,1\}$ let $R_{n,i}$ be the stored reward and $K_{n,i}$ the nonnegative discounted live-state kernel. A common law parameter $\lambda$ is used at all dates:
\begin{equation}
 R_n^\lambda=(1-\lambda)R_{n,0}+\lambda R_{n,1},\qquad
 K_n^\lambda=(1-\lambda)K_{n,0}+\lambda K_{n,1}.
 \label{eq:mixture}
\end{equation}
The current action is chosen before the endpoint-law draw. Discharge rewards are included once in $R$; lost live-state mass is not assigned an additional terminal reward. Later-date controls consist of a union of two fixed meshes, with 1,565 common actions, three frozen date- and state-dependent proposals, and surrender when permitted. Consumption is in $[0.05,0.8]$, deliberate drift in $[-0.2,0.2]$, and common financial shares in $[-0.5,0.8]$. In the no-adjustment regime, only zero-drift actions are admissible; stochastic preference variation remains.

At the first date, consumption and adjustment pairs are finite but the risky share is optimized over the stored piecewise-affine interpolation of its breakpoints. The positive mandate requires $0<\pi_0\leq L$ throughout the first interval; the other requires $-S\leq\pi_0\leq0$. The application fixes $L=0.8$ and $S=0.5$. Subsequent shares are not restricted by the initial sign. A closure value at $\pi_0=0$ supplies an upper supremum for the positive class, not an attained positive policy. Section~\ref{sec:arithmetic} gives the corresponding lower-witness construction.

Elective surrender at a nonboundary live state pays $G-F$ and is available only at dates $n\geq m$. Forced discharge pays $G$ without the fee. The contract $m=8$ has no elective surrender within the operating horizon. We compare contracts at $\lambda=0.125$ and $d=0.42425$. Baseline purchaser primitives are $b=1$, $C(E)=0.02E^2$, $D(m)=0.02(m-1)/8$, $\rho_F=\nu=0$, $\alpha_C=\zeta=1$, and $G_0=G(2,1.25)$. These are normalized institutional parameters, not estimates from a dataset. Their alternatives are evaluated in Section~\ref{sec:results}.

\section{Response Sets and Procurement}
\label{sec:response}
\subsection{A sharp finite-query characterization}
Suppose values are enclosed at $t_0=0,t_1,\ldots,t_M$:
\begin{equation}
 L_i\leq W(t_i)\leq U_i,\qquad i=0,\ldots,M.
 \label{eq:queries}
\end{equation}
An $\eta$-response satisfies $J(p;0)\geq W(0)-\eta$. Exact responses correspond to $\eta=0$ when the supremum is attained. Define a lifted set with variables $v,a,z$ and one pair $(\beta_i,z_i)$ for each query. Here $v$ represents the center value, $a$ the candidate's center payoff, and $(\beta_i,z_i)$ an affine response supporting query $i$. Impose
\begin{align}
 &L_0\leq v\leq U_0,\quad v-\eta\leq a\leq v,\quad z,z_i\in K,
 \label{eq:oracle1}\\
 &a+t_j\cdot z\leq U_j,\qquad \beta_i\leq v,
 \label{eq:oracle2}\\
 &\beta_i+t_j\cdot z_i\leq U_j,\quad
   \beta_i+t_i\cdot z_i\geq L_i,\quad \beta_0=v.
 \label{eq:oracle3}
\end{align}
Let $\mathcal Z_\eta(L,U)$ be its projection onto $z$. All the restrictions are linear when $K$ is a polytope. Compact feature bounds also give finite intercept bounds: $\beta_i\geq L_i-h_K(t_i)$, where $h_K$ is the support function of $K$.

\begin{theorem}[Finite-query response identification]
\label{thm:oracle}
Assume $K$ is a nonempty compact polytope, the query intervals are mutually consistent with an affine-response economy, and each queried supremum is attained in that economy or its stated closure. Then every attained $\eta$-response has feature vector in $\mathcal Z_\eta(L,U)$. Conversely, each feasible lifted point defines a finite affine-response economy satisfying all the query intervals and having the candidate as an $\eta$-response. Hence support bounds over $\mathcal Z_\eta(L,U)$ are sharp relative to the query information and $K$.
\end{theorem}

The necessity argument uses the candidate's affine payoff and an optimal response at each query. For sufficiency, take those same affine functions as the entire response menu. Their maximum at zero is $v$; each query is at least its designated lower bound and no affine function exceeds its upper bound. The construction proves sharpness without claiming that every feasible lifted point is implementable in the already specified settlement economy. That distinction allows the theorem to be used as a valid outer bound there while reserving the equality-graph restrictions for stronger fixed-economy statements.

One-dimensional secants are a special case. With duration in $[0,1]$, queries at $d-h,d,d+h$ give
\begin{equation}
 \max\!\left\{0,\frac{L(d)-U(d-h)-\eta}{h}\right\}
 \leq A(p)\leq
 \min\!\left\{1,\frac{U(d+h)-L(d)+\eta}{h}\right\}.
 \label{eq:secants}
\end{equation}
Decreasing $h$ reduces the span of economic curvature but magnifies value error and $\eta$. There is no assertion that an arbitrarily small secant is better. With two features, separate bounds generally discard their dependence. For instance, at $W(t)=\max\{t_1,t_2\}$, coordinate queries at $\pm h$ do not rule out a response with both feature coordinates equal to one. Adding $(h,h)$ rules it out and bounds their sum by one in the exact oracle. The supplement gives the finite-menu witnesses and rational certificates for this example and its interval variants.

\subsection{Checking support bounds rather than trusting an optimizer}
Write a bounded linear program as $\max\{c^\top x:Ax\leq b,\ \ell\leq x\leq u\}$. For any nonnegative vector $\mu$, weak duality gives
\begin{equation}
 c^\top x\leq \mu^\top b+
 \sum_j\max\{(c-A^\top\mu)_j\ell_j,(c-A^\top\mu)_ju_j\}.
 \label{eq:dual}
\end{equation}
The right side remains valid if $\mu$ is not exactly dual feasible in the usual stationarity sense: the residual is charged through the box support. In the implementation, a floating-point optimizer proposes $\mu$, while rational arithmetic checks nonnegativity, every coefficient, and the complete residual term. A solver's success flag is therefore neither a feasibility proof nor a certificate of the reported upper bound. The same argument applied to $-c$ gives a lower support bound.

\subsection{From response identification to institutional choice}
For each contract, let $[L_j,U_j]$ enclose private value and let $\mathcal Z_j$ contain all $\eta$-response moments. Define $s_j(z)=bA+\rho_F FH+\nu Q$ and
\begin{align}
 \underline q_j&=[G_0-U_j+\alpha_C C(E)]_+,\qquad
 \overline q_j=[G_0-L_j+\eta+\alpha_C C(E)]_+,\label{eq:grants}\\
 \underline\Pi_j&=\inf_{z\in\mathcal Z_j}s_j(z)-\overline q_j-(1-\alpha_C)C(E)-D(m),\label{eq:lowerprofit}\\
 \overline\Pi_j&=\sup_{z\in\mathcal Z_j}s_j(z)-\underline q_j-(1-\alpha_C)C(E)-D(m).\label{eq:upperprofit}
\end{align}
The grant $\overline q_j$ induces participation for every feasible $\eta$-response. The upper bound deliberately allows a rival its most favorable response and its smallest possibly sufficient grant, even when these are not jointly attainable. Retaining their dependence in a lifted program can sharpen this upper bound but is not necessary for its validity.

\begin{theorem}[Robust procurement]
\label{thm:procurement}
Suppose the value and response enclosures above are valid, and the candidate admits an attained exact response when $\eta=0$, or feasible $\eta$-responses when $\eta>0$. If
\begin{equation}
 \underline\Pi_{j_*}>\max\{0,\sup_{j\ne j_*}\overline\Pi_j\},
 \label{eq:separation}
\end{equation}
then offering $j_*$ with grant $\overline q_{j_*}$ strictly dominates nonprocurement and every other participating offer, for every admitted response of the candidate and every admitted response of every rival. More generally, if the right side is at most $\underline\Pi_{j_*}+\varepsilon$, the candidate is executable with regret at most $\varepsilon$ relative to the most favorable competing response.
\end{theorem}

This is a result about contracts with endogenous response, not merely about the output of a policy optimizer. Exact behavior is the economic benchmark. The $\eta$ statement measures robustness to approximate optimization, and is not obtained by substituting an exact-policy support calculation into an $\eta$-response claim. An outside option of zero for the purchaser is explicitly included.

\begin{proposition}[Efficient capacity and incidence]
\label{prop:capacity}
Suppose capacity affects neither operating payoffs nor feasible policies except through \eqref{eq:enforcement}, and $C$ is nondecreasing. For every feasible $(F,E)$ there is a weakly better purchaser contract with the same $F$ and $E=F/\zeta$. This holds with binding or slack participation. Under binding participation, changing $\alpha_C$ leaves purchaser surplus unchanged. Changing $\rho_F$ instead adds $\Delta\rho_F\,FH$ and generally changes the ranking.
\end{proposition}

Indeed, total capacity and participation expenditure is
$[G_0-W+\alpha_C C]_++(1-\alpha_C)C$, a nondecreasing function of $C$. Under binding participation it equals $G_0-W+C$. Thus the equality of fee and capacity in the baseline can be an optimum of an explicit enforcement technology rather than a restriction imposed on two otherwise independent instruments. This proposition does not apply to a capacity payment that reduces managed wealth. That intervention changes $W$ and is not an incidence relabeling.

\section{Continuous Enforcement without an Unchecked Menu Approximation}
\label{sec:continuous}
By Proposition~\ref{prop:capacity}, the baseline continuous instrument is $F\in[0,1]$, with $E=F$. Fix a term and mandate, and write a policy's payoff at the central benefit as $B-FH$, absorbing $dA$ into $B$. A value query at $(d+\delta,F_i)$ implies
\begin{equation}
 B+\delta A-F_iH\leq U(d+\delta,F_i).
 \label{eq:continuousquery}
\end{equation}
These inequalities constrain the \emph{same} policy moments across fee cells. At an attained exact optimum at fee $F_i$, let $H_i\in[\underline H_i,\overline H_i]$. Its affine payoff supplies a lower bound at another fee:
\begin{equation}
 W(d,F)\geq L(d,F_i)-(F-F_i)h_i,
 \quad h_i=\begin{cases}\overline H_i,&F\geq F_i,\\
 \underline H_i,&F\leq F_i.\end{cases}
 \label{eq:feesupport}
\end{equation}
It follows that an $\eta$-response at $F$ must satisfy $B-FH\geq L(d,F_i)-(F-F_i)h_i-\eta$.

For a cell $F\in[a,b]$, comparison with an exact response at a strictly exterior fee gives further monotonicity bounds:
\begin{equation}
 H\leq\overline H_i+\frac{\eta}{a-F_i}\quad(F_i<a),\qquad
 H\geq\underline H_i-\frac{\eta}{F_i-b}\quad(F_i>b).
 \label{eq:feerange}
\end{equation}
For exact responses, endpoint bounds apply by the same two-policy argument. The inequalities follow from adding the optimality inequalities at the two fees; no derivative of $W$ is taken. Intersect them with $[0,1]$ to obtain $H\in[\underline h,\overline h]$.

Introduce $T=FH$. Its graph over this rectangle is contained in the four affine inequalities
\begin{align}
 T&\geq aH+\underline hF-a\underline h,&
 T&\geq bH+\overline hF-b\overline h,\nonumber\\
 T&\leq bH+\underline hF-b\underline h,&
 T&\leq aH+\overline hF-a\overline h.
 \label{eq:product}
\end{align}
A supporting affine function of $C(F/\zeta)$ gives a lower bound on capacity cost throughout the cell. With $C(F)=cF^2$ and midpoint $v$, it is $c(2vF-v^2)$. Combining \eqref{eq:continuousquery}--\eqref{eq:product}, participation, and moment bounds yields a linear program whose objective is the purchaser's payoff. The upper program may admit a fictitious response or a relaxed product. That enlarges its feasible set and preserves the direction of the bound.

\begin{theorem}[A continuum procurement certificate]
\label{thm:continuous}
For each term and mandate, let finitely many closed fee cells cover the entire feasible interval. On every cell impose valid query, response, product, capacity-cost, and participation relaxations as above. Let $U_I$ be a rationally verified upper objective and let $L_*$ be the lower payoff of an executable candidate. Then
\begin{equation}
 \sup_{j,p,q\;\mathrm{participating}}\Pi_j(p,q)-L_*
 \leq\max\{0,\max_I U_I\}-L_*.
 \label{eq:continuumregret}
\end{equation}
In particular, if the candidate lower payoff exceeds the maximum upper bound on every cell with a different term, its compulsory term is identified even if the optimal fee is not attained or the remaining same-term regret is positive.
\end{theorem}

The last qualification matters at a surrender indifference. The least fee eliminating surrender may leave a low-service best response at the boundary. A slightly larger fee can implement high service, while the infimum cost itself is not achieved under adverse tie selection. Our reported object is consequently an executable contract with a bound against the global supremum, not an asserted maximizer at an untested kink. Every retained cell appears in the witness; a collection of favorable evaluation points cannot pass as a continuum certificate.

The algorithm begins with the deposited fee queries, adds benefit perturbations where the response upper bound is loose, and subdivides the cell with the largest verified purchaser upper bound. It stops at a specified regret tolerance or an explicit computational budget. No convergence flag is interpreted as economic separation. A reported tolerance is met only when the final cover and its inequalities establish \eqref{eq:continuumregret}. Wider benefit perturbations can be valuable because they constrain implausible high-duration responses that remain compatible with very local value queries. The choice of query is thus directed by the purchaser's unresolved comparison.

\section{Contract Choice in the Settlement Economy}
\label{sec:results}
\subsection{Nested menus and continuous fees}
Table~\ref{tab:menus} reports exact-response comparisons on three nested menus. The six-point menu and the two refinements use the same 1,617-state economy, the same later-date actions, the same first-date interpolation, and the same public permissions. Only the set of purchasable fees changes. All ties within a contract are enclosed by an independently reconstructed Bellman equality graph, rather than resolved in favor of the purchaser. The displayed rival is the contract furnishing the largest competing upper payoff.
\begin{table}[t]
\caption{Nested Enforcement Menus}\label{tab:menus}
\centering
\footnotesize
\begin{tabular}{@{}rlcrcr@{}}
\toprule
Step & Regime & $(F,m,\sigma)$ & Surplus lower & Rival & Margin \\
\midrule
0.20 & Adjustment & $0.60,3,+$ & 0.0960214 & $0.80,1,+$ & 0.00059979 \\
0.20 & No adjustment & $0.60,4,+$ & 0.0580640 & $0.80,1,-$ & 0.00022405 \\
0.10 & Adjustment & $0.60,3,+$ & 0.0960214 & $0.70,2,+$ & 0.00009979 \\
0.10 & No adjustment & $0.30,6,+$ & 0.0583927 & $0.60,4,+$ & 0.00032847 \\
0.05 & Adjustment & $0.65,2,+$ & 0.0972714 & $0.55,3,+$ & 0.00009979 \\
0.05 & No adjustment & $0.55,4,+$ & 0.0590227 & $0.75,2,+$ & 0.00000842 \\
\bottomrule
\end{tabular}
\par\smallskip\begin{minipage}{\linewidth}\footnotesize Target: canonical manifest \texttt{54adb353c5b85210d} (full hash in the release manifest). Exact responses; value allowance $10^{-7}$ and support allowance $10^{-9}$. Lower bounds and margins are rounded down. Each margin compares the candidate lower bound against every other offer upper bound and nonprocurement.\end{minipage}
\end{table}


The original three-versus-four term comparison is therefore correct for its institution but not invariant to menu refinement. This is economically informative: a fee slightly above a surrender threshold can substitute for a compulsory interval. A coarser menu makes that substitution more expensive because it requires excess capacity. A claim about commitment length that omits the enforceability menu can conflate this discretization effect with the preference-adjustment mechanism.

The continuous experiment applies Theorem~\ref{thm:continuous} on $F\in[0,1]$ after separately optimizing capacity under \eqref{eq:enforcement}. Its final candidates, global upper bounds, and term-specific competing bounds appear in Table~\ref{tab:continuous}. Both identified compulsory terms are one interval. The adjusted contract requires a smaller charge than the no-adjustment contract, so the continuously optimized comparison concerns the price of enforcement rather than the earlier difference in term length. These calculations establish the stated regret guarantees on the continuous fee interval; they do not approximate continuous consumption or adjustment choices. The latter remain part of the unchanged settlement target.
\begin{table}[t]
\caption{Continuous-Fee Procurement}\label{tab:continuous}
\centering
\footnotesize
\begin{tabular}{@{}lrrrrr@{}}
\toprule
Regime & Fee & Term & Surplus lower & Global upper & Regret upper \\
\midrule
Adjustment & 0.72880859 & 1 & 0.097598279 & 0.097599163 & 0.000000883 \\
No adjustment & 0.80087891 & 1 & 0.059936055 & 0.059936149 & 0.000000093 \\
\bottomrule
\end{tabular}
\par\smallskip\begin{minipage}{\linewidth}\footnotesize Target: canonical manifest \texttt{54adb353c5b85210d} (full hash in the release manifest). Both mandates, all eight terms, and every $F\in[0,1]$ are covered; $E=F$ is derived from the capacity theorem. Behavioral tolerance is $\eta=0$. Regret and upper bounds are rounded up. Exact input fees and outward grants: Adjustment: $F=0.72880859375$, $q=0.792634995286$; No adjustment: $F=0.80087890625$, $q=0.830297219351$.\end{minipage}
\end{table}


The numerical comparisons have distinct logical forms. A positive finite-menu margin excludes every other menu item. A continuous-fee certificate bounds a supremum over a covered interval and can identify a term without identifying a unique fee. An exact-response equality graph is not automatically an $\eta$-response graph. For the original menu we retain the separate secant certificate with $\eta=10^{-8}$ and the service-price and term-cost rectangle. Its smallest deposited margin is $7.6880857\times10^{-5}$. For the continuous comparison, the behavioral tolerance is reported in the table, not inherited implicitly from the original menu.

\subsection{Joint service and surrender information}
For a fee-receiving purchaser the relevant operating output is $A+FH$. Table~\ref{tab:joint} fixes actual settlement contracts near surrender transitions, evaluates coordinate reward queries, and then adds the two joint directions. In each case the lifted response set uses the same feature box and value-error allowance. The upper bound falls strictly; it is not merely a smaller error bar on a separately estimated derivative. The gain arises because a single policy cannot simultaneously select the most favorable duration from one response and the most favorable fee receipt from another without satisfying the joint value inequalities.
\begin{table}[t]
\caption{Joint Queries Bound Duration and Fee Receipts}\label{tab:joint}
\centering
\footnotesize
\begin{tabular}{@{}lcrrrr@{}}
\toprule
Regime & Mandate & $\eta$ & Coordinate upper & Joint upper & Reduction \\
\midrule
Adjustment & + & 0.0000 & 0.955820 & 0.882345 & 0.073475 \\
Adjustment & + & 0.0001 & 0.963945 & 0.887345 & 0.076600 \\
Adjustment & $-$ & 0.0000 & 0.915092 & 0.882145 & 0.032947 \\
Adjustment & $-$ & 0.0001 & 0.923217 & 0.887145 & 0.036072 \\
No adjustment & + & 0.0000 & 0.891508 & 0.890449 & 0.001059 \\
No adjustment & + & 0.0001 & 0.899133 & 0.895449 & 0.003684 \\
No adjustment & $-$ & 0.0000 & 0.888871 & 0.888088 & 0.000783 \\
No adjustment & $-$ & 0.0001 & 0.896496 & 0.893088 & 0.003408 \\
\bottomrule
\end{tabular}
\par\smallskip\begin{minipage}{\linewidth}\footnotesize Target: canonical manifest \texttt{54adb353c5b85210d} (full hash in the release manifest). The objective is $A+FH$. Adjustment uses $(m,F)=(2,0.625)$ and no adjustment $(4,0.525)$. The coordinate oracle uses $h=0.02$ and five queries; the joint oracle adds $(h,Fh)$ and its negative. Query error is $10^{-10}$. Bounds concern the same policy feature vector, not independent coordinate selections.\end{minipage}
\end{table}


This experiment gives the finite-query theorem a direct economic use beyond the one-dimensional secant argument. It does not claim that the oracle set equals the feasible response set of the settlement dynamic program. Exact equality-graph supports can be tighter when the full Bellman solution is available. Finite-query supports are useful when value certificates, rather than complete feasible-response descriptions, are the transferable objects, and when the purchaser changes the weights on several service features.

\subsection{Institutional sensitivity}
Table~\ref{tab:institutions} records counterfactuals on the 0.05 fee menu. Each row uses actual settlement response moments. Capacity cost is evaluated at the efficient level $E=F/\zeta$; changes in fee recipient use joint duration--surrender supports; quality changes use the wealth-qualified duration moment. These are not substitutions from the separate analytical contract example in the supplement.
\begin{table}[t]
\caption{Institutional Counterfactuals in the Same Settlement Economy}\label{tab:institutions}
\centering
\footnotesize
\begin{tabular}{@{}lcrcr@{}}
\toprule
Institution & Adjusted offer & Margin & Fixed offer & Margin \\
\midrule
Baseline & $0.65,2,+$ & 0.0000997 & $0.55,4,+$ & 0.0000084 \\
All fees to purchaser & $0.65,2,+$ & 0.0000997 & $0.60,3,+$ & 0.0002224 \\
Half fees to purchaser & $0.65,2,+$ & 0.0000997 & $0.55,4,+$ & 0.0001202 \\
Capacity paid by purchaser & $0.65,2,+$ & 0.0000997 & $0.55,4,+$ & 0.0000084 \\
Capacity cost shared & $0.65,2,+$ & 0.0000997 & $0.55,4,+$ & 0.0000084 \\
Efficiency $0.8$ & $0.35,5,+$ $\dagger$ & -0.0000003 & $0.00,7,+$ & 0.0001960 \\
Efficiency $1.2$ & $0.75,1,+$ & 0.0005553 & $0.85,1,+$ & 0.0002775 \\
Capacity coefficient $0.01$ & $0.75,1,+$ & 0.0007747 & $0.85,1,+$ & 0.0008747 \\
Capacity coefficient $0.04$ & $0.35,5,+$ & 0.0000998 & $0.00,7,+$ & 0.0007947 \\
Cubic capacity cost & $0.65,2,+$ & 0.0003347 & $0.75,2,+$ & 0.0000811 \\
Quadratic term cost & $0.45,4,+$ & 0.0004372 & $0.55,4,+$ & 0.0000210 \\
Quality premium $+0.05$ & $0.65,2,+$ & 0.0000997 & $0.55,4,+$ & 0.0000082 \\
Quality premium $-0.05$ & $0.65,2,+$ & 0.0000997 & $0.80,1,-$ & 0.0006062 \\
Fees and quality premium & $0.65,2,+$ & 0.0000997 & $0.60,3,+$ & 0.0002050 \\
Outside value $+0.05$ & $0.65,2,+$ & 0.0000997 & $0.55,4,+$ & 0.0000084 \\
Outside value $-0.05$ & $0.65,2,+$ & 0.0000997 & $0.55,4,+$ & 0.0000084 \\
Service price $0.9$ & $0.65,2,+$ & 0.0000997 & $0.00,1,+$ & 0.0004559 \\
Service price $1.1$ & $0.65,2,+$ & 0.0000997 & $0.75,2,+$ & 0.0000110 \\
\bottomrule
\end{tabular}
\par\smallskip\begin{minipage}{\linewidth}\footnotesize Target: canonical manifest \texttt{54adb353c5b85210d} (full hash in the release manifest). All rows use the 0.05 fee menu and exact-response graph supports. Offers are $(F,m,\sigma)$. A dagger denotes an unresolved strict ordering, not a certified winner. Capacity is $F/\zeta$, the charge ceiling remains one, and participation is evaluated with its positive part. These are finite-menu sensitivities; they are not continuous-fee claims.\end{minipage}
\end{table}


Three distinctions organize the results. First, shifting a numeraire capacity payment between parties leaves the allocation unchanged when participation binds, as Proposition~\ref{prop:capacity} predicts. The outside-value shifts in the table change required compensation but not the operating incentives. Second, transferring surrender fees to the purchaser changes the social value of a response with elective stopping and can change the chosen term. This is not a transfer-incidence invariance. Third, the shape of enforcement and term costs changes the tradeoff between monetary and temporal commitment. Curvature and enforcement efficiency can eliminate the baseline term difference even while adjustment changes private continuation values.

Quality valuation is equally substantive. Rewarding duration at wealth at least $1.25$ is not the same service as undifferentiated operating time. A negative premium can change the preferred initial financial mandate in the no-adjustment regime. The service definition is therefore not innocuous. The method handles the alternative by keeping the joint moment, not by declaring all service measures equivalent. A signing grant entering wealth, a different discharge rule, or a different action set requires new dynamic values and is outside these incidence-only calculations.

The remaining sensitivity concerns the numerical query scale. For the original winning offers, the deposited $h$ checks at $0.0025$, $0.005$, and $0.01$ retain positive no-adjustment margins of approximately $3.04\times10^{-5}$, $1.14\times10^{-4}$, and $1.56\times10^{-4}$ under the specified conservative construction. These are secant/error tradeoffs, not proof that the continuous model has negligible approximation error. The complete queries, opponents, and error conventions are supplied with the tables.

\section{Adjustment Opportunities and Financial Exposure}
\label{sec:mechanism}
\subsection{An exact decomposition}
Let $V^a$ and $V^0$ denote continuation values with and without deliberate adjustment, and let $O=V^a-V^0$. At an adjusted optimal first action in class $\sigma$, keep consumption and investment fixed and replace its drift by zero. Suppose this counterfactual action is feasible. Let $K_\sigma$ be its discounted continuation row. Define $g_\sigma$ as the gain from replacing that zero-drift action by the adjusted action evaluated with $V^a$, and $B_\sigma$ as the zero-drift counterfactual's payoff evaluated with $V^0$ minus the optimal no-adjustment payoff. Then $B_\sigma\leq0$ and
\begin{equation}
 \Omega_\sigma=g_\sigma+K_\sigma O+B_\sigma,\qquad
 \Omega_+-\Omega_-=g_+-g_-+(K_+-K_-)O+B_+-B_-.
 \label{eq:exposure}
\end{equation}
This identity neither asserts that the signed exposure term is positive nor identifies a primitive economic ordering. For $D=K_+-K_-=D^+-D^-$ and an enclosure $\ell\leq O\leq u$, its sharp box bound is
\begin{equation}
 D^+\ell-D^-u\leq DO\leq D^+u-D^-\ell.
 \label{eq:signedexposure}
\end{equation}
Taking an absolute norm before this multiplication loses the direction of exposure. Likewise, showing that negative drift is dominated does not order the two financial kernels or the distribution of future opportunities.

\subsection{Primitive conditions in an ordered stopping subclass}
A useful sufficient condition can be stated for a restricted but explicit economy. Let wealth follow an exogenous order-preserving Markov transition, with common state-independent survival and discount factors. Before an irreversible preference upgrade, the agent may choose its exercise date. Consumption is an exogenous nondecreasing function of wealth taking values in $(0,1]$. An upgrade raises the preference index from $u$ to $u'>u$ and entails a wealth-independent cost. The terminal increment from upgrading is nondecreasing in wealth (it may be zero), and the upgrade does not change the wealth transition or consumption rule. There is no additional wealth-dependent surrender option in this subclass. The incremental flow is
\begin{equation}
 \Delta f(c)=\frac{c^{1-u'}}{1-u'}-\frac{c^{1-u}}{1-u}.
\end{equation}
Although its level can have either sign, its derivative is $c^{-u'}-c^{-u}\geq0$. Coupled ordered wealth paths therefore order the present value of exercising at any common stopping time.

\begin{proposition}[Ordered adjustment opportunities]
\label{prop:primitive}
Under the preceding primitives, the value of the irreversible upgrade option is nondecreasing in wealth at every date. If the positive initial financial transition first-order stochastically dominates the nonpositive transition with the same total discounted mass, then its future opportunity exposure is weakly larger. With equal current gains and equal replacement terms, \eqref{eq:exposure} gives $\Omega_+\geq\Omega_-$. Strict inequality follows if the stochastic dominance is strict on a set where the opportunity value is strictly increasing and that set receives positive exposure.
\end{proposition}

The assumptions locate the mechanism rather than hiding it in a value-function ordering. Monotone consumption, a single upgrade, common survival, and the absence of wealth-dependent surrender are doing work. The full settlement economy has endogenous consumption, continuous preference adjustment on its mesh, and elective surrender. It need not satisfy this proposition. The decomposition remains valid there, but its sign must be checked rather than imported from the subclass.

\subsection{A map including adverse cases}
The settlement map uses $\lambda\in\{0,0.125,1\}$, $d\in\{0.35,0.42425,0.5\}$, $F\in\{0,0.4,0.8\}$, and $m\in\{1,4\}$. All 54 points are retained. Each entry recomputes the adjusted and no-adjustment problems, the zero-drift counterfactual rows, the signed exposure, and wealth-bin contributions. Table~\ref{tab:mechanism} summarizes the sign frequencies and ranges; the complete rows make the reversals visible.
\begin{table}[t]
\caption{Adjustment-Opportunity Map}\label{tab:mechanism}
\centering
\footnotesize
\begin{tabular}{@{}lrrrrr@{}}
\toprule
Component & Positive & Negative & Near zero & Minimum & Maximum \\
\midrule
Current-gain difference & 16 & 11 & 27 & -0.0019671 & 0.0000544 \\
Future signed exposure & 15 & 12 & 27 & -0.0036131 & 0.0003893 \\
Financial replacement & 0 & 0 & 54 & 0.0000000 & 0.0000000 \\
Total option difference & 15 & 12 & 27 & -0.0055802 & 0.0004436 \\
\bottomrule
\end{tabular}
\par\smallskip\begin{minipage}{\linewidth}\footnotesize Target: canonical manifest \texttt{54adb353c5b85210d} (full hash in the release manifest). All 54 combinations of three laws, three benefits, three fees, and two terms are retained. Counts use a reporting tolerance of $2\times10^{-9}$; full signed intervals and values are deposited. The map is not a proof of a sign between grid points.\end{minipage}
\end{table}


The map shows why a local sign is not a transferable primitive theorem. Changing the fee or the time at which surrender is permitted changes the opportunity process reached by each initial financial action. The current drift gain can be positive while the future signed contribution reverses the class difference. The replacement term is zero in the recorded map; this is an outcome of these first-action choices, not a general cancellation imposed on the decomposition. The primitive proposition and the adverse map serve complementary roles: one identifies sufficient structure, and the other records where the unrestricted application departs from it.

\section{Neural Bellman Operators and the Computational Comparison}
\label{sec:computation}
\subsection{Proposal, evaluation, improvement, and certification}
For a discounted finite-horizon problem, let $T_n$ be the Bellman operator and $T_n^p$ its fixed-policy counterpart. A neural operator proposes a value or policy indexed by economic parameters. Policy evaluation yields an executable lower payoff; Bellman improvement searches the admissible actions; an independent upper construction bounds what that search or proposal might have missed. These are separate mathematical objects. A proposed policy with evaluated value $L$ and a valid optimal upper bound $U$ has regret at most $U-L$, regardless of how it was proposed.

The compendium proves the continuous operator, approximation, recursive-utility, and equilibrium versions with their own regularity and contraction assumptions. Here the changing-law certificate uses the mixture structure in \eqref{eq:mixture}. A fixed policy's value is a Bernstein polynomial in the common law parameter. Its coefficients can be evaluated from the deposited policy and kernels. Restricting the polynomial to a subinterval by de Casteljau recursion gives a lower witness throughout that interval. An upper relaxation conditions on counts of endpoint draws while preserving a common current action before the next draw. Signed chord corrections provide a compressed alternative. The data-dependent arithmetic account applies separately to policy coefficients, upper recursions, restrictions, and signed comparisons.

These constructions do not assume that the common law parameter can be selected afresh at every state. A rectangular parameter-lifting bound permits such extra choices and is consequently an informative, potentially looser comparator. Neither signed compression nor neural approximation is guaranteed to dominate an exact finite-state calculation. Their costs and certificate widths must be compared on a common target with the same lower witnesses and the same error convention.

\subsection{Matched evidence and its limits}
Table~\ref{tab:proposals} includes teacher construction, fitting, online evaluation, and the final policy gap for the actual deposited neural, polynomial, and nearest-anchor proposals. The training laws are $0.12$ and $0.13$; tests include $0$, $0.125$, $0.25$, $0.75$, and $1$. The network is a trained two-hidden-layer, 32-unit model, not a relabeled interpolation scheme. The target is the same 1,617-state, eight-interval settlement economy used by the economic results.
\begin{table}[t]
\caption{Proposal Costs and Verified Policy Gaps}\label{tab:proposals}
\centering
\footnotesize
\begin{tabular}{@{}llrrrrr@{}}
\toprule
Regime & Proposal & Teacher & Fit & Online max. & Gap max. & KiB \\
\midrule
Adjustment & Neural & 1.576 & 1.125 & 1.742 & $0.806$ & 22.1 \\
Adjustment & Polynomial & 1.576 & 0.005 & 1.762 & $0.0902$ & 5.8 \\
Adjustment & Nearest & 1.576 & 0.000 & 1.710 & $2.13\times10^{-07}$ & 101.1 \\
No adjustment & Neural & 1.798 & 0.750 & 1.897 & $0.568$ & 13.4 \\
No adjustment & Polynomial & 1.798 & 0.003 & 1.870 & $0.0147$ & 1.8 \\
No adjustment & Nearest & 1.798 & 0.000 & 1.795 & $2\times10^{-07}$ & 101.1 \\
\bottomrule
\end{tabular}
\par\smallskip\begin{minipage}{\linewidth}\footnotesize Target: canonical manifest \texttt{54adb353c5b85210d} (full hash in the release manifest). Seconds are from the completed R13 execution, not a timing rerun in R14. Online work includes proposal, policy evaluation, exact-DP upper reference, and independent verification. Maxima cover five test laws and both mandates. KiB reports stored proposal parameters or anchor policies, not process peak memory. No fitting step is needed for nearest-anchor lookup; teacher cost remains charged. The experiment uses one target size and one recorded neural seed.\end{minipage}
\end{table}


At this scale the nearest-anchor proposal is the strongest of the three in achieved policy value. The network's larger gaps are charged rather than excused by its fit. Online comparison includes an exact upper solve and policy verification, which limits the benefit from a cheap proposal. The measurements consequently do not establish a neural advantage in state dimension, action complexity, or end-to-end accuracy per resource. The title denotes the maintained operator architecture and its theory; the numerical procurement theorem relies on verified response and value inequalities, not on a claim of neural superiority.

The supplemental computational table reports common-law chord and count bounds over $[0.12,0.13]$, $[0,0.25]$, and $[0,1]$, together with the independently specified rectangular-lifting implementation. Endpoint solves and lower-witness work are charged to every method using them. The comparison with parameter lifting is a comparison of algorithms on the same arrays, not an external-package speed ranking. Peak memory and hardware identities are attached to their execution records. Repeated timing distributions, high-dimensional scaling, and failure probabilities across random training runs are not inferred from a single recorded fit. A missing measurement is marked as such rather than entered as zero.

\subsection{Arithmetic and an open financial mandate}
\label{sec:arithmetic}
All binary64 entries of the canonical target are treated as exact rational inputs to its mathematical Bellman problem. Evaluation error is then derived under a stated arithmetic model: correctly rounded basic operations, round-to-nearest, finite nonnegative sparse data, no overflow, gradual underflow, and reductions within the documented operation-count envelope. The derivation is not a formal proof of every possible implementation of a sparse numerical library.

Let $s$ bound nonzero terms per sparse row, $M$ discounted row mass, and $B_{n+1}$ a bound on the magnitude of continuation values. With unit roundoff $u$ and $\gamma_k=ku/(1-ku)$, a row and reward calculation has a forward-error envelope of the form
\begin{equation}
 e_n\leq M e_{n+1}+\gamma_{k_n}\bigl(\bar R_n+M B_{n+1}\bigr)+k_n\,\mathrm{sub},
 \label{eq:roundoff}
\end{equation}
where $k_n$ also covers mixture, reward, and comparison operations and $\mathrm{sub}$ charges subnormal absolute error. Maximization is nonexpansive in the sup norm. The actual audit derives $s$, $M$, $B_n$, and each operation envelope from all stored rows. Its largest bound across the inherited certificate operations is below $7.057\times10^{-11}$. The inherited economic records charge $10^{-7}$ per value; new value queries charge $10^{-10}$ only after the audit verifies that this allowance covers their operations. Equality-graph support calculations charge their separately derived support allowance.

For a positive first mandate, let $\bar W_+$ be the optimum on $[0,L]$. If its optimizing financial knot is strictly positive, the associated policy is an attained feasible witness. If it is zero, choose a positive share $\epsilon$ in the adjacent affine segment while keeping the same first consumption--adjustment pair and continuation policy. The value loss is the segment slope times $\epsilon$ and is charged explicitly. This produces a feasible lower payoff, while $\bar W_+$ remains an upper supremum. An exact-response procurement statement may use the open class only when attainment has been established; otherwise it uses feasible near-optimal responses and the charged closure gap. The equality of an open supremum and its closure value never makes the zero action feasible.

\section{From a Numerical Target to an Economic Decision}
\label{sec:transfer}
\subsection{A complete conditional error chain}
Let $T_n$ denote the Bellman operator of an intended model and $\widehat T_n$ that of a numerical target. An interpolation or lifting operator $I_n$ places numerical functions on the model's state space. Suppose, for a stated envelope of admissible continuation functions,
\begin{equation}
 \|T_n I_{n+1}v-I_n\widehat T_n v\|\leq\delta_n,
 \qquad \|I_n v-I_n w\|\leq\Lambda_n\|v-w\|.
 \label{eq:residual}
\end{equation}
The norm, domain, boundary treatment, and continuation envelope are part of this assumption. A residual evaluated only at selected nodes does not establish it. Write the local allowance as
\begin{equation}
 \delta_n\leq\delta_n^{\rm time}+\delta_n^{\rm trunc}
 +\delta_n^{\rm state}+\delta_n^{\rm action}
 +\delta_n^{\rm first}+\delta_n^{\rm kernel}
 +\delta_n^{\rm boundary}+\delta_n^{\rm arithmetic},
 \label{eq:budget}
\end{equation}
and let $\delta_N^{\rm terminal}$ bound the terminal discrepancy. Each entry refers to a specified intermediate operator, so errors are not omitted or counted twice. In particular, first-date financial interpolation and later-date consumption--adjustment discretization are separate entries.

\begin{theorem}[Model-to-decision transfer]
\label{thm:transfer}
Suppose the model Bellman operators are Lipschitz in continuation value with factors $M_n$, the lifting and residual bounds \eqref{eq:residual} hold, and terminal discrepancy is at most $\delta_N$. Then the value discrepancy is at most $\Delta_n$, where
\begin{equation}
 \Delta_N=\delta_N,\qquad \Delta_n=\delta_n+M_n\Delta_{n+1}.
 \label{eq:transferrecursion}
\end{equation}
At each required reward query enlarge $[L_i,U_i]$ to $[L_i-\Delta_i,U_i+\Delta_i]$. If the same affine response features and feasible-policy interpretation hold in the intended model, Theorems~\ref{thm:oracle} and~\ref{thm:procurement} apply to these enlarged intervals. If policies are transferred instead, a uniform two-sided policy-payoff discrepancy $\Delta$ maps an intended-model $\eta$-response to a target $(\eta+2\Delta)$-response, provided the stipulated policy correspondence is available. A service-moment discrepancy must be charged separately when the transferred feature itself changes.
\end{theorem}

The first route avoids a policy-selection error by querying values in the intended model with valid intervals. The second is convenient for policy witnesses but requires a two-sided correspondence; transferring only a favorable numerical policy does not bound all intended-model competitors. For a one-dimensional secant, enlarging both value endpoints by $\Delta$ increases each service error by at most $2\Delta/h$. Participation adds at most $\Delta$ per contract. Hence a simple two-contract sufficient decision-error charge is $2(2b\Delta/h+\Delta)$, before quality, fee, or feature-transfer errors. The smallest purchaser gap, not the magnitude of private values, determines whether this budget is adequate.

\subsection{What the deposited constructor establishes}
The canonical arrays are produced from stochastic-settlement primitives by a finite-step transition construction and state projection. They are not exact transition probabilities of the original continuous diffusion. The binary64 audit starts after that construction. To examine the missing approximation step directly, the new diagnostic selects the interior first action $(c,\theta,\pi)=(0.05,0,0)$ at $u=2$ and computes the normalized stored preference variance with exact rational arithmetic on its row. Both endpoint laws give approximately
\begin{equation}
 \operatorname{Var}_{\rm stored}(u_1)=0.0008838834764831865,
 \qquad \sigma_u^2\Delta t=0.0003125.
 \label{eq:variance}
\end{equation}
For a drift-free increment smaller than a preference grid cell, linear projection to the adjacent nodes contributes variance $\Delta u\,\sigma_u\sqrt{\Delta t}$, rather than $\sigma_u^2\Delta t$. Here their ratio is $2.828427\ldots$. This is a local moment comparison, not a lower or upper bound on value error. It nevertheless rules out identifying array-evaluation accuracy with constructor accuracy.

A raw total-variation comparison between an atomless diffusion transition and a grid-supported law is also inappropriate: they are mutually singular before a common projection or lifting is specified. The relevant operator estimate must use an explicit interpolation and a continuation regularity class, or compare consistently projected kernels. The supplement states the corresponding Lipschitz and smooth-test-function alternatives and keeps the boundary and stopping terms separate.

The inherited stress calculation supplies a sufficient neighborhood under hypothetical reward and kernel errors. It does not establish that the diffusion lies in that neighborhood. Accordingly, the paper's executed strict inequalities are about the hash-identified settlement array economy and the continuous \emph{contract} extension within it. The general model-to-decision theorem specifies what is required for the antecedent continuous-state model; its quantitative inclusion has not been established by the deposited constructor. This distinction preserves, rather than replaces, the continuous-state theory and provides a testable approximation requirement for it.

\section{Conclusion}
A dynamic value certificate becomes an economic contract certificate only after the induced response, participation payment, and enforcement cost are treated together. Finite reward queries can identify a joint outer set of all responses without differentiability or a favorable tie-breaking rule. Rational support bounds make that information usable by a purchaser. A complete cover of the enforcement interval then replaces an unchecked menu approximation with a bound against every continuous instrument.

The settlement application illustrates why each step matters. Refining enforcement changes the compulsory-term comparison in the original institution; fee incidence and quality valuation alter which response the purchaser values; and signed future exposure can reverse a local adjustment advantage. These findings do not disappear when a neural proposal is inaccurate, because proposal quality and certificate validity are separate. Nor does a verified finite-array calculation by itself establish a diffusion approximation. The contribution is the explicit chain from dynamic value information to response restrictions and then to an executable institutional choice, with each numerical and modeling error charged at the comparison where it matters.
